\section{Conclusiones}
\label{sec:conclusiones}

Al terminar el diseño, la programación, las pruebas y la evaluación, se llegó a las siguientes conclusiones, en orden con los objetivos:

\begin{itemize}
    \item \textbf{Cumplimiento del Objetivo General:} Se hizo y se puso a funcionar el SISTEMA WEB DE INSCRIPCIÓN PARA EL MÓDULO TECNOLÓGICO PRODUCTIVO SAN JULIÁN BTH --- VERSIÓN 2.0, hecho en tres partes (lo que se ve, lo que organiza y donde se guardan los datos), logrando una página completa que bajó el registro de varios días de fila a 3 a 5 minutos por postulante.
    \item \textbf{Diagnóstico de Deficiencias Operativas (OE1):} Se encontraron los tres problemas del proceso 2022 a 2025 (contraseñas sin recuperación, papeles sin registro en sistema y paralelos armados a mano) con preguntas en el primer bimestre de 2026 a 60 estudiantes postulantes y 5 funcionarios, ordenados en la Tabla~\ref{tab:brechas_diagnostico}.
    \item \textbf{Recuperación de contraseña:} La recuperación sola con carnet y fecha de nacimiento resolvió el problema más grande de la gestión pasada, dejando que los postulantes recuperen su acceso en menos de 2 minutos sin ir a ventanilla y sin perder su cupo, según la Tabla~\ref{tab:comparacion_antes_despues}.
    \item \textbf{Papeles de matrícula en orden:} Se unió la preinscripción con la matrícula con un espacio que registra y revisa en el sistema los papeles que pide el BTH (libretas de 1ro a 3ro, cédula de tutores, carnet y certificado), logrando todos los expedientes en orden y sin papeles perdidos, conforme a la Tabla~\ref{tab:comparacion_antes_despues}.
    \item \textbf{Reparto en cursos y paralelos:} Se acabaron los choques de horario entre el MÓDULO TECNOLÓGICO PRODUCTIVO SAN JULIÁN y los colegios del distrito, porque el sistema reparte solo a los aceptados en su curso (4to A, 4to B) según su colegio, dejando los talleres equilibrados.
    \item \textbf{Pruebas con muchos usos a la vez:} Las pruebas con más de 100 usos por minuto mostraron que el sistema aguanta, responde en menos de 1.2 segundos y mantiene todos los cupos en orden en las siete especialidades, sin duplicar ningún cupo.
    \item \textbf{Facilidad de uso y calidad:} Las preguntas a 60 estudiantes con puntaje de 1 a 5 mostraron una facilidad de uso de 92 \% (al menos 90 \%) y que los errores bajaron de 35 \% a menos de 5 \%, gracias a la revisión automática al escribir.
    \item \textbf{Ahorro para las familias:} Se acabaron las filas de noche a la intemperie de las familias de San Julián, ahorrando más de 20000 Bs por gestión en pasajes y copias, con una inversión de 13574 Bs, gasto de 190 Bs/mes y menos de 0.50 Bs por postulante, recuperándose en menos de una gestión.
\end{itemize}

En consecuencia, lo planteado en la Sección~\ref{sec:objetivos} se cumplió: la Versión 2.0 con recuperación por carnet y fecha de nacimiento y guardado seguro que no duplica cupos bajó el tiempo en más de 90 \%, resolvió la contraseña en menos de 2 minutos, dejó papeles y cupos en orden al 100 \% sin duplicar, como se vio en las pruebas con muchos usos a la vez, en la revisión de datos y en la Tabla~\ref{tab:comparacion_antes_despues} del Capítulo~7.
